\documentclass[12pt,a4paper]{article}
% Drake_preamble.tex - LaTeX Preamble Template
% 作者: 謝慕揚, MD, PhD, FESC
% 
% 使用說明:
% 1. 表格請使用 longtable，不要有垂直分隔線 |
% 2. 表格內換行使用 \newline，不要使用 \\
% 3. 藥物名稱、檢驗、檢查請維持英文
% 4. Reference 格式請使用 \href 加上驗證過的 DOI

% Including essential packages for Chinese typesetting
\usepackage{xeCJK}
\usepackage{fontspec}
% Configuring Chinese font with Noto Serif CJK TC, fallback to SimSun
\IfFontExistsTF{Noto Serif CJK TC}{
  \setCJKmainfont{Noto Serif CJK TC}
  \setCJKsansfont{Noto Serif CJK TC}
  \setCJKmonofont{Noto Serif CJK TC}
}{\setCJKmainfont{SimSun}
  \setCJKsansfont{SimSun}
  \setCJKmonofont{SimSun}
}

% Including other necessary packages
\usepackage{geometry}
\usepackage{titlesec}
\usepackage{enumitem}
\usepackage{xcolor}
\usepackage{colortbl}
\usepackage{tcolorbox}
\usepackage{booktabs}
\usepackage{longtable}
\usepackage{array}
\usepackage{graphicx}
\usepackage{tikz}
\usepackage{fancyhdr}
\usepackage{multicol}
\usepackage{datetime2}
\usepackage{hyperref}
\usepackage{mdframed}
\usepackage{amsmath}
\usepackage{amssymb}
\usepackage{multirow}

\hypersetup{
    colorlinks=true,
    linkcolor=emergency_blue,
    citecolor=emergency_blue,
    urlcolor=emergency_blue,
    pdfborder={0 0 0}
}

% Defining page geometry
\geometry{
  top=2.5cm,
  bottom=2.5cm,
  left=2cm,
  right=2cm,
  headheight=25pt
}

% Adjusting topmargin to compensate for increased headheight
\addtolength{\topmargin}{-10pt}

% Defining colors
\definecolor{emergency_red}{RGB}{186,24,27}
\definecolor{emergency_blue}{RGB}{0,114,188}
\definecolor{emergency_gray}{RGB}{120,120,120}
\definecolor{highlight_yellow}{RGB}{255,250,205}

% Setting title formats
\titleformat{\section}[block]{\Large\bfseries\color{emergency_red}}{\thesection}{10pt}{}
\titleformat{\subsection}[block]{\large\bfseries\color{emergency_blue}}{\thesubsection}{8pt}{}
\titleformat{\subsubsection}[block]{\normalsize\bfseries\color{emergency_gray}}{\thesubsubsection}{6pt}{}

% Defining custom environment for important notes
\newmdenv[
    linecolor=emergency_red,
    backgroundcolor=highlight_yellow,
    linewidth=2pt,
    topline=true,
    bottomline=true,
    leftline=true,
    rightline=true,
    innertopmargin=10pt,
    innerbottommargin=10pt,
    innerleftmargin=10pt,
    innerrightmargin=10pt
]{importantbox}

% Defining note command with tcolorbox
\newcommand{\note}[1]{
    \par
    \noindent
    \begin{tcolorbox}[
        colback=emergency_blue!5!white,
        colframe=emergency_blue,
        title=\textbf{重點提醒},
        fonttitle=\bfseries,
        boxrule=0.5pt,
        arc=3pt,
        left=6pt,
        right=6pt,
        top=6pt,
        bottom=6pt
    ]
    #1
    \end{tcolorbox}
}

% Key findings box
\newcommand{\keyfinding}[1]{
    \par
    \noindent
    \begin{tcolorbox}[
        colback=yellow!10!white,
        colframe=orange!75!black,
        title=\textbf{核心發現摘要},
        fonttitle=\bfseries,
        boxrule=0.5pt,
        arc=3pt
    ]
    #1
    \end{tcolorbox}
}

% Defining custom date format for ISO 8601
\DTMnewdatestyle{iso}{
  \renewcommand{\DTMdisplaydate}[4]{\number##1-\number##2-\number##3}
  \renewcommand{\DTMDisplaydate}[4]{\number##1-\number##2-\number##3}
}
\DTMsetdatestyle{iso}
\newcommand{\mydate}{\DTMtoday}


% Custom header for this document
\pagestyle{fancy}
\fancyhf{}
\fancyhead[L]{\textcolor{emergency_red}{JACC 教學講義}}
\fancyhead[R]{\textcolor{emergency_blue}{Smartwatch AF Detection}}
\fancyfoot[C]{\textcolor{emergency_gray}{\thepage}}
\renewcommand{\headrulewidth}{1pt}
\renewcommand{\headrule}{\hbox to\headwidth{\color{emergency_red}\leaders\hrule height \headrulewidth\hfill}}

\begin{document}

%============================================================
% Title Page
%============================================================
\begin{center}
{\LARGE\bfseries\textcolor{emergency_red}{Journal of the American College of Cardiology}}\\[0.5cm]
{\Large\textbf{Enhanced Detection and Prompt Diagnosis of Atrial Fibrillation Using Apple Watch: The EQUAL Trial}}\\[0.3cm]
{\large 使用 Apple Watch 增強心房顫動之偵測與即時診斷：EQUAL 隨機對照試驗}\\[0.5cm]
{\normalsize \href{https://www.jacc.org}{JACC 2026; Epub ahead of print}}\\[0.3cm]
{\normalsize 整理：謝慕揚 MD, PhD, FESC \quad 日期：\mydate}
\end{center}

\vspace{0.5cm}

%============================================================
\section{研究背景}
%============================================================

心房顫動 (atrial fibrillation, AF) 是最常見的心律不整，與中風風險顯著相關。消費型智慧穿戴裝置（如 Apple Watch）已具備偵測 AF 的功能數年，但過去缺乏將此技術整合至臨床工作流程的隨機對照試驗證據。

\subsection{研究目的}
評估智慧手錶結合遠距醫療工作流程，是否能提升高中風風險族群的 AF 偵測率。

%============================================================
\section{研究方法}
%============================================================

\begin{itemize}
    \item \textbf{研究設計}：隨機對照試驗 (Randomized Controlled Trial)
    \item \textbf{研究機構}：Amsterdam University Medical Center 與 Cardiology Centers of the Netherlands
    \item \textbf{受試者}：220 位高 AF 風險且高中風風險的患者
    \item \textbf{介入措施}：
    \begin{itemize}
        \item 介入組：配戴 Apple Watch 六個月
        \item 對照組：標準照護（未配戴智慧手錶）
    \end{itemize}
    \item \textbf{工作流程整合}：患者記錄的 ECG 由 E-health 團隊於 24 小時內判讀
\end{itemize}

%============================================================
\section{主要發現}
%============================================================

\keyfinding{
\begin{itemize}
    \item 介入組的 AF 診斷率為對照組的 \textbf{4 倍}（fourfold increase）
    \item 介入組的 AF 診斷時間較對照組更早
    \item 24 小時內 ECG 判讀的臨床工作流程運作良好
\end{itemize}
}

\subsection{ECG 記錄觸發機制}
患者在以下情況記錄 ECG：
\begin{enumerate}
    \item 出現症狀時（symptomatic）
    \item 收到裝置的不規則心律通知（irregular rhythm notification）
\end{enumerate}

%============================================================
\section{臨床工作流程}
%============================================================

\begin{importantbox}
\textbf{Heartwacht 系統整合}

荷蘭心臟科中心原本使用手持式行動 ECG 裝置（medical device）進行遠距監測。本研究將 Apple Watch 整合至此既有系統：

\begin{enumerate}
    \item 患者透過 Apple Watch 記錄 ECG
    \item ECG 資料傳送至 Heartwacht 平台
    \item E-health 團隊於 24 小時內判讀
    \item 若偵測到 AF，立即啟動後續處置
\end{enumerate}
\end{importantbox}

%============================================================
\section{臨床意涵}
%============================================================

\subsection{適用族群}
\begin{itemize}
    \item 年齡 $\geq$ 65 歲
    \item 高 AF 風險
    \item 高 AF 相關中風風險（高 CHA$_2$DS$_2$-VASc 分數）
\end{itemize}

\subsection{實施考量}

\begin{longtable}{p{4cm}p{10cm}}
\toprule
\textbf{面向} & \textbf{說明} \\
\midrule
\endfirsthead
\toprule
\textbf{面向} & \textbf{說明} \\
\midrule
\endhead
\bottomrule
\endfoot

工作流程整合 & 需建立專責團隊判讀 ECG，否則消費型穿戴裝置的 ECG 難以有效利用 \\
\midrule
給付機制 & 荷蘭已有既存的遠距監測給付架構，新技術可順利整合\newline 其他國家需建立相應的給付機制 \\
\midrule
地理優勢 & 遠距監測對於幅員廣大、就醫距離遠的國家尤其重要 \\
\bottomrule
\end{longtable}

%============================================================
\section{未來研究方向}
%============================================================

\subsection{REGAL Trial}
目前進行中的大型試驗，研究目標包括：
\begin{itemize}
    \item 硬性臨床終點（hard clinical endpoints）
    \item 中風風險
    \item 認知功能退化（cognitive decline）
\end{itemize}

\note{
\textbf{Screen-detected AF 的治療適應症}

目前仍需更多證據確認：篩檢偵測到的 AF（screen-detected AF）在高 CHA$_2$DS$_2$-VASc 分數患者中，是否應常規給予抗凝血治療。
}

\subsection{其他心律不整偵測}
研究團隊正探索智慧手錶偵測其他心律不整的潛力：
\begin{itemize}
    \item 高度房室傳導阻滯（high-grade AV block）
    \item 傳導異常（conduction disorders）
    \item 其他間歇性心律不整（intermittent arrhythmias）
\end{itemize}

%============================================================
\section{重點摘要}
%============================================================

\begin{enumerate}
    \item \textbf{顯著效益}：智慧手錶整合遠距監測可提升 4 倍 AF 偵測率
    \item \textbf{關鍵成功因素}：24 小時內 ECG 專業判讀的臨床工作流程
    \item \textbf{適用對象}：65 歲以上、高 AF 及中風風險患者
    \item \textbf{待解答問題}：Screen-detected AF 的最佳治療策略
    \item \textbf{未來應用}：可能擴展至其他心律不整的偵測
\end{enumerate}

%============================================================
\section{參考文獻}
%============================================================

\begin{enumerate}
    \item Van Steijn NJ, Blommestijn IS, Blok S, et al. Enhanced detection and prompt diagnosis of atrial fibrillation using Apple Watch: a randomized controlled trial. \textit{J Am Coll Cardiol}. 2026; Epub ahead of print.

    \item \href{https://www.tctmd.com/news/smartwatch-increases-af-diagnosis-older-high-risk-patients-equal}{TCTMD News: Smartwatch Increases AF Diagnosis in Older, High-risk Patients: EQUAL. January 2026.}

    \item Mannhart D, Lopes Ventura Ferreira V, Baur J, et al. Clinical Validation of 5 Direct-to-Consumer Wearable Smart Devices to Detect Atrial Fibrillation: BASEL Wearable Study. \href{https://doi.org/10.1016/j.jacep.2022.09.011}{\textit{JACC Clin Electrophysiol}. 2023;9(2):232-242.}
\end{enumerate}

\end{document}
